\begin{lstlisting}[style=pythonstyle, caption={app.py}, label={lst:app.py}]
import transcribe
import CreateComicPage
import TextToScript
import scriptToImage
import scriptToPrompt

def main():
    transcribe.main()
    TextToScript.main()
    scriptToPrompt.main()
    scriptToImage.main()
    CreateComicPage.main()

if __name__ == "__main__":
    main()
\end{lstlisting}

\begin{lstlisting}[style=pythonstyle, caption={transcribe.py}, label={lst:transcribe.py}]
import whisper
import torch
from transformers import AutoTokenizer, AutoModelForTokenClassification
from pyannote.audio import Pipeline
import json
import re

# Load Whisper model
whisper_model = whisper.load_model("turbo")

# Load punctuation restoration model
tokenizer = AutoTokenizer.from_pretrained("oliverguhr/fullstop-punctuation-multilang-large")
punct_model = AutoModelForTokenClassification.from_pretrained("oliverguhr/fullstop-punctuation-multilang-large")
punct_model.eval()

# Load pyannote diarization pipeline (this requires your HuggingFace token)
diarization_pipeline = Pipeline.from_pretrained("pyannote/speaker-diarization-3.1")
diarization_pipeline.to(torch.device("cuda"))


def transcribe_with_segments(audio_path):
    """Transcribe audio and return timestamped segments."""
    result = whisper_model.transcribe(audio_path)
    return result.get("segments", [])

def restore_punctuation(text):
    """Restore punctuation and capitalization using the punctuation model."""
    inputs = tokenizer(text, return_tensors="pt")
    with torch.no_grad():
        logits = punct_model(**inputs).logits
    predictions = torch.argmax(logits, dim=-1)[0].tolist()
    tokens = tokenizer.convert_ids_to_tokens(inputs['input_ids'][0])

    # Reconstruct text properly handling SentencePiece tokens
    words = []
    current_word = ""
    
    for i, (token, label) in enumerate(zip(tokens, predictions)):
        # Skip special tokens
        if token in ['<s>', '</s>', '<pad>', '<unk>']:
            continue
            
        # Handle SentencePiece tokens (they start with ▁ which is \u2581)
        if token.startswith('▁') or token.startswith('\u2581'):
            # Start of new word
            if current_word:
                words.append(current_word)
            current_word = token.replace('▁', '').replace('\u2581', '')
        elif token.startswith('##'):
            # BERT-style subword continuation
            current_word += token[2:]
        else:
            # Regular token or continuation
            if current_word and not token.startswith('▁') and not token.startswith('\u2581'):
                current_word += token
            else:
                if current_word:
                    words.append(current_word)
                current_word = token.replace('▁', '').replace('\u2581', '')
        
        # Add punctuation based on label
        if label == 1:  # PERIOD
            current_word += "."
        elif label == 2:  # COMMA
            current_word += ","
        elif label == 3:  # QUESTION_MARK
            current_word += "?"
    
    # Don't forget the last word
    if current_word:
        words.append(current_word)
    
    # Join words and clean up
    punctuated_text = " ".join(words)
    
    # Clean up any remaining special characters and normalize spaces
    punctuated_text = re.sub(r'[▁\u2581]', '', punctuated_text)
    punctuated_text = re.sub(r'\s+', ' ', punctuated_text)
    punctuated_text = punctuated_text.strip()
    
    # Capitalize first letter
    if punctuated_text:
        punctuated_text = punctuated_text[0].upper() + punctuated_text[1:]
    
    return punctuated_text

def diarize_audio(audio_path):
    """Run diarization and return speaker segments with timestamps."""
    diarization = diarization_pipeline(audio_path)
    diarized_segments = []
    for turn, _, speaker in diarization.itertracks(yield_label=True):
        diarized_segments.append({
            "start": turn.start,
            "end": turn.end,
            "speaker": speaker
        })
    return diarized_segments

def merge_transcript_and_diarization(transcript_segments, diarized_segments):
    """Match each transcript segment to a speaker based on overlap."""
    merged = []

    for t_seg in transcript_segments:
        # Find diarization segments overlapping this transcript segment
        candidates = [d for d in diarized_segments if not (d["end"] < t_seg["start"] or d["start"] > t_seg["end"])]

        if candidates:
            # Pick the speaker with the longest overlap
            best_speaker = None
            max_overlap = 0
            for c in candidates:
                overlap_start = max(t_seg["start"], c["start"])
                overlap_end = min(t_seg["end"], c["end"])
                overlap = overlap_end - overlap_start
                if overlap > max_overlap:
                    max_overlap = overlap
                    best_speaker = c["speaker"]
        else:
            best_speaker = "Unknown"

        merged.append({
            "start": t_seg["start"],
            "end": t_seg["end"],
            "speaker": best_speaker,
            "text": t_seg["text"]
        })

    return merged

def create_clean_transcript_with_speakers(audio_path, output_json):
    print("Transcribing audio with Whisper...")
    transcript_segments = transcribe_with_segments(audio_path)

    print("Restoring punctuation...")
    for seg in transcript_segments:
        seg["text"] = restore_punctuation(seg["text"])

    print("Running speaker diarization...")
    diarized_segments = diarize_audio(audio_path)

    print("Merging transcript and diarization...")
    merged = merge_transcript_and_diarization(transcript_segments, diarized_segments)

    with open(output_json, "w", encoding="utf-8") as f:
        json.dump(merged, f, indent=2, ensure_ascii=False)

    print(f"Transcript with speakers saved to {output_json}")

def main():
    audio_file = "../audio/fullSession2.wav"
    output_file = "output/clean_transcript_with_speakers.json"
    create_clean_transcript_with_speakers(audio_file, output_file)

if __name__ == "__main__":
    main()

\end{lstlisting}

\begin{lstlisting}[style=pythonstyle, caption={TextToScript.py}, label={lst:TextToScript.py}]
import ollama
import json
import re
import time

def load_transcript(path):
    with open(path, "r", encoding="utf-8") as f:
        return json.load(f)

def adaptive_panel_generator(transcript, chunk_size=15, max_retries=3):
    panels = []
    position = 0
    memory_buffer = []
    story_so_far = ""

    while position < len(transcript):
        chunk = transcript[position : position + chunk_size]
        print(f"Processing chunk at line {position}...")

        retries = 0
        while retries < max_retries:
            try:
                used_indices, panel, summary = generate_panel_from_chunk(chunk, story_so_far)
                if panel:
                    panels.append(panel)

                    memory_buffer.append(summary)
                    if len(memory_buffer) > 10:
                        memory_buffer.pop(0)
                    story_so_far = "\n".join(memory_buffer)

                    num_used = max(used_indices) + 1 if used_indices else 1
                    position += num_used
                    break  # success, break out of retry loop
                else:
                    retries += 1
                    print(f"No panel returned, retrying... ({retries}/{max_retries})")
                    time.sleep(0.3)
            except Exception as e:
                retries += 1
                print(f"Panel generation error: {e} (retry {retries}/{max_retries})")
                time.sleep(0.3)

        else:
            print("❌Failed after maximum retries, skipping chunk.")
            position += 1  # give up and move on

    return panels

def generate_panel_from_chunk(chunk, story_so_far):
    content = "\n".join(f"{i}: {seg['speaker']}: {seg['text']}" for i, seg in enumerate(chunk))

    prompt = f"""
You are an AI comic scriptwriter adapting a Dungeons & Dragons session.

Here is what has happened so far:
{story_so_far.strip()}

Here is the next part of the transcript:
{content}

IMPORTANT FILTERING:

- Ignore any rule talk, dice rolls, or out-of-character chatter.
- Focus only on in-character dialogue, world-building, and story-driving actions.

TASK:

Generate ONE new panel based on a selection of lines from this chunk. Your job is to condense meaningful story progression into a single visual scene.

Output JSON in this format:
{{
  "used_segments": [list of indices],
  "panel": {{
    "setting": "...",
    "characters": [Character names],
    "objects": [Important objects],
    "action": "...",
    "mood": "...",
    "camera_view": "...",
    "lighting": "...",
    "text": "..."
  }},
  "summary": "One sentence summary of the panel for story memory."
}}

- Every panel must include all 9 fields and cannot be empty.
- Output only valid JSON — no extra text, no commentary.
"""

    response = ollama.chat(
        model="llama3.1",
        messages=[{"role": "user", "content": prompt}],
        stream=False,
    )

    output = response["message"]["content"]
    return parse_adaptive_json(output)

def parse_adaptive_json(text):
    try:
        match = re.search(r'{.*}', text, re.DOTALL)
        if not match:
            raise ValueError("No JSON found")
        json_text = match.group(0)
        json_text = re.sub(r",\s*}", "}", json_text)
        json_text = re.sub(r",\s*\]", "]", json_text)
        parsed = json.loads(json_text)

        used = parsed.get("used_segments", [])
        panel = parsed.get("panel", {})
        summary = parsed.get("summary", "")

        return used, panel, summary
    except Exception as e:
        print("JSON parsing failed:", e)
        return [], None, ""

def save_panels(panels, output_path):
    with open(output_path, "w", encoding="utf-8") as f:
        json.dump(panels, f, indent=2, ensure_ascii=False)

def main():
    transcript_path = "output/clean_transcript_with_speakers.json"
    output_path = "output/adaptive_comic_script.json"

    transcript = load_transcript(transcript_path)
    panels = adaptive_panel_generator(transcript)
    save_panels(panels, output_path)
    print(f"✅ Script saved to {output_path}")

if __name__ == "__main__":
    main()

\end{lstlisting}

\begin{lstlisting}[style=pythonstyle, caption={scriptToPrompt.py}, label={lst:scriptToPrompt.py}]
import json
import ollama
import time

def normalize_panel(panel):
    def normalize_entry(entry):
        if isinstance(entry, str):
            return entry
        elif isinstance(entry, dict):
            name = entry.get("name", "Unknown")
            extras = [v for k, v in entry.items() if k != "name" and v]
            return f"{name} ({', '.join(extras)})" if extras else name
        return str(entry)

    panel["characters"] = [normalize_entry(c) for c in panel.get("characters", [])]
    panel["objects"] = [normalize_entry(o) for o in panel.get("objects", [])]

    return panel

def generate_prompt_with_llama(panel):
    system_message = "You are a prompt engineer for Stable Diffusion XL (SDXL)."

    user_message = f"""
Create a visually descriptive SDXL prompt based on the following panel description:

Setting: {panel.get('setting', '')}
Characters: {", ".join(panel.get("characters", []))}
Objects: {", ".join(panel.get("objects", []))}
Action: {panel.get('action', '')}
Mood: {panel.get('mood', '')}
Camera View: {panel.get('camera_view', '')}
Lighting: {panel.get('lighting', '')}

Requirements:
- Combine all of the above into a single natural-language description.
- Maximum length: 75 tokens.
- Do NOT add quotes or commentary, just return the prompt text only.
- The prompts need to be short, concise, and suitable for generating a single comic book panel.
"""

    response = ollama.chat(
        model="llama3.1",
        messages=[
            {"role": "system", "content": system_message},
            {"role": "user", "content": user_message}
        ],
        stream=False
    )

    prompt = response["message"]["content"].strip()
    return prompt

def process_script(input_path, output_path):
    with open(input_path, "r", encoding="utf-8") as f:
        script = json.load(f)

    for i, panel in enumerate(script):
        print(f"Generating SDXL prompt for panel {i}...")
        try:
            panel = normalize_panel(panel)
            prompt = generate_prompt_with_llama(panel)
            panel["sdxl_prompt"] = prompt
            panel["negative_prompt"] = (
                "deformed, blurry, extra limbs, poorly drawn hands, extra fingers, "
                "mutated, out of frame, ugly, tiling, low resolution, bad anatomy, watermark"
            )
        except Exception as e:
            print(f"Failed for panel {i}: {e}")
            panel["sdxl_prompt"] = ""
            panel["negative_prompt"] = ""
        time.sleep(0.3)

    with open(output_path, "w", encoding="utf-8") as f:
        json.dump(script, f, indent=2, ensure_ascii=False)

    print(f"SDXL prompts saved to {output_path}")

def main():
    process_script(
        input_path="output/adaptive_comic_script.json",
        output_path="output/adaptive_comic_script_with_prompts.json"
    )

if __name__ == "__main__":
    main()

\end{lstlisting}

\begin{lstlisting}[style=pythonstyle, caption={scriptToImage.py}, label={lst:scriptToImage.py}]
from diffusers import DiffusionPipeline
import torch
import AddTextToImage
import json

testing = False

def generate_images_from_script(script_path):
    # Load comic panel script with LLM-generated prompts
    with open(script_path, "r", encoding="utf-8") as f:
        comic_script = json.load(f)

    # Load SDXL pipeline
    pipe = DiffusionPipeline.from_pretrained(
        "stabilityai/stable-diffusion-xl-base-1.0",
        torch_dtype=torch.float16,
        use_safetensors=True,
        variant="fp16"
    )
    pipe.load_lora_weights("../lora/Fantasy_art_XL_V1.safetensors")
    pipe.to("cuda")

    for i, panel in enumerate(comic_script):
        prompt = "Comic book illustration. Vivid fantasy." + panel.get("sdxl_prompt", "")
        negative_prompt = panel.get("negative_prompt", "")

        if not prompt:
            print(f"Skipping panel {i}: no prompt found.")
            continue

        image = pipe(
            prompt=prompt,
            negative_prompt=negative_prompt,
        ).images[0]

        image.save(f"output/images/imagesWithoutText/panel_{i}.png")
        image_with_text = AddTextToImage.add_text_to_panel(panel["text"], image)
        image_with_text.save(f"output/images/imagesWithText/panel_{i}.png")

        if i == 5 and testing:
            return

def main():
    generate_images_from_script("output/adaptive_comic_script_with_prompts.json")

if __name__ == "__main__":
    main()

\end{lstlisting}

\begin{lstlisting}[style=pythonstyle, caption={CreateComicPage.py}, label={lst:CreateComicPage.py}]
from PIL import Image
import os

def resize_and_add_border(image, target_size, border_size):
    resized_image = Image.new("RGB", target_size, "black")
    resized_image.paste(image, ((target_size[0] - image.width) // 2, (target_size[1] - image.height) // 2))
    return resized_image

def load_images():
    number_of_images = 6
    images_loaded = []
    image_files = os.listdir("output/images/imagesWithText")

    # Sort filenames by panel index (assumes filenames like panel_0.png)
    sorted_files = sorted(
        image_files,
        key=lambda name: int("".join(filter(str.isdigit, name)))
    )

    for i in range(len(sorted_files)):
        image_path = os.path.join("output/images/imagesWithText", sorted_files[i])
        image = Image.open(image_path)
        images_loaded.append(image)

    return images_loaded


def main():
    images = load_images()
    columns, rows = 2, 3
    images_per_page = columns * rows

    for i in range(0, len(images), images_per_page):
        page_images = images[i:i + images_per_page]

        # Calculate the size of the output image
        output_width = columns * page_images[0].width + (columns - 1) * 10
        output_height = rows * page_images[0].height + (rows - 1) * 10

        result_image = Image.new("RGB", (output_width, output_height), "black")

        # Combine images into a grid with black borders
        for j, img in enumerate(page_images):
            x = (j % columns) * (img.width + 10)
            y = (j // columns) * (img.height + 10)

            resized_img = resize_and_add_border(img, (page_images[0].width, page_images[0].height), 10)
            result_image.paste(resized_img, (x, y))

        result_image = result_image.resize((1024, 1536))
        result_image.save(f"comicPages/comicPage_{i // images_per_page}.jpg")
        print(f"Saved comic page: comicPage_{i // images_per_page}.jpg")

if __name__ == "__main__":
    main()
\end{lstlisting}

\begin{lstlisting}[style=pythonstyle, caption={pdfMerger.py}, label={lst:pdfMerger.py}]
from PIL import Image
import os

image_dir = "comicPages"
image_files = [f for f in os.listdir(image_dir) if f.lower().endswith((".jpg", ".jpeg", ".png"))]
image_files.sort()

images = []
for file in image_files:
    img_path = os.path.join(image_dir, file)
    img = Image.open(img_path).convert("RGB")
    images.append(img)

if images:
    images[0].save("comic.pdf", save_all=True, append_images=images[1:])
\end{lstlisting}